\exercisesfor{4}

The conventions of § 4 remain valid unless otherwise mentioned.

\begin{enumerate}
\def\labelenumi{\arabic{enumi}.}
\item
  Let \(\mathfrak{g}\) be a nilpotent Lie algebra and \(p\) (resp. \(q\) ) the smallest integer such that \(\mathcal{C}^p\mathfrak{g} = \{0\}\) (resp. \(\mathcal{C}_q\mathfrak{g} = \mathfrak{g}\) ). Show that \(p = q + 1\) and that \(\mathcal{C}_i\mathfrak{g} \supset \mathcal{C}^{p - i}\mathfrak{g}\) . (Use the argument of Proposition 1.)
\item
  Let \(\mathfrak{g}\) be a semi-direct product of an algebra \(\mathfrak{h}\) of dimension 1 and a commutative ideal \(\mathfrak{g}'\) . Let \(x \in \mathfrak{h}, x \neq 0\) , and \(u\) be the restriction of \(\operatorname{ad}_{\mathfrak{g}} x\) to \(\mathfrak{g}'\) .
\end{enumerate}

\begin{enumerate}
\def\labelenumi{(\alph{enumi})}
\item
  For g to be nilpotent, it is necessary and sufficient that u be nilpotent.
\item
  For the Killing form of \(\mathfrak{g}\) to be zero, it is necessary and sufficient that \(\operatorname{Tr}(u^2) = 0\) .
\item
  Deduce from (a) and (b) that there exist non-nilpotent Lie algebras whose Killing forms are zero.
\item
  Deduce from (a) that in a nilpotent Lie algebra such that \(\mathcal{C}^{p-1}\mathfrak{g} \neq \{0\}\) , \(\mathcal{C}^p\mathfrak{g} = \{0\}\) it is possible that \(\mathcal{C}_i\mathfrak{g} \neq \mathcal{C}^{p-i}\mathfrak{g}\) .
\end{enumerate}

\begin{enumerate}
\def\labelenumi{\arabic{enumi}.}
\setcounter{enumi}{2}
\tightlist
\item
  \begin{enumerate}
  \def\labelenumii{(\alph{enumii})}
  \tightlist
  \item
    Let \(\mathfrak{g}\) be a nilpotent Lie algebra, \(\mathfrak{z}\) its centre and \(\mathfrak{b}\) a non-zero ideal of \(\mathfrak{g}\) . Show that \(\mathfrak{z} \cap \mathfrak{b} \neq \{0\}\) . (Consider \(\mathfrak{b}\) as a \(\mathfrak{g}\) -module by means of the adjoint representation.)
  \end{enumerate}
\end{enumerate}

\begin{enumerate}
\def\labelenumi{(\alph{enumi})}
\setcounter{enumi}{1}
\tightlist
\item
  If in a Lie algebra \(\mathfrak{g}\) an ideal \(\mathfrak{b}\) is contained in \(\mathcal{C}_{i+1}\mathfrak{g}\) but not in \(\mathcal{C}_i\mathfrak{g}\) , show that the ideals \(\mathfrak{b} \cap \mathcal{C}_k\mathfrak{g}\) are all distinct for \(0 \leqslant k \leqslant i + 1\) . (Apply (a).)
\end{enumerate}

\begin{enumerate}
\def\labelenumi{\arabic{enumi}.}
\setcounter{enumi}{3}
\tightlist
\item
  Let g be a nilpotent Lie algebra.
\end{enumerate}

\begin{enumerate}
\def\labelenumi{(\alph{enumi})}
\item
  Every subalgebra of \(\mathfrak{g}\) is subinvariant. (Use Proposition 3.)
\item
  Let \(\mathfrak{b}\) be a vector subspace of \(\mathfrak{g}\) such that \(\mathfrak{b} + \mathcal{D}\mathfrak{g} = \mathfrak{g}\) . Show that the subalgebra of \(\mathfrak{g}\) generated by \(\mathfrak{b}\) is \(\mathfrak{g}\) . (Apply (a) to this subalgebra. Reduce it thus to the case where \(\mathfrak{b}\) is an ideal of \(\mathfrak{g}\) and use Proposition 4 of § 1.) Deduce that the minimum number of generators of \(\mathfrak{g}\) is \(\dim \mathfrak{g}/\mathcal{D}\mathfrak{g}\) .
\end{enumerate}

\begin{enumerate}
\def\labelenumi{\arabic{enumi}.}
\setcounter{enumi}{4}
\tightlist
\item
  \begin{enumerate}
  \def\labelenumii{(\alph{enumii})}
  \tightlist
  \item
    Show that a Lie algebra g in which every subalgebra is subinvariant is nilpotent. (Show that if \(\dim g > 1\) every \(x \in g\) belongs to an ideal \(h \neq g\) of g which has the same property as g and hence is nilpotent by induction on the dimension of g; hence \(ad_{g} x\) is nilpotent.)
  \end{enumerate}
\end{enumerate}

\begin{enumerate}
\def\labelenumi{(\alph{enumi})}
\setcounter{enumi}{1}
\tightlist
\item
  Show that if in a Lie algebra \(\mathfrak{g}\) every subalgebra distinct from \(\mathfrak{g}\) is distinct from its normalizer, \(\mathfrak{g}\) is nilpotent. (Reduce it to (a).)
\end{enumerate}

\begin{enumerate}
\def\labelenumi{\arabic{enumi}.}
\setcounter{enumi}{5}
\item
  Let g be a nilpotent Lie algebra and a a commutative ideal of g. The following conditions are equivalent: (a) a is a maximal commutative ideal of g; (b) a is a maximal commutative subalgebra of g: (c) a is equal to its centralizer \(a'\) in g. (The implications not (a) \(\Rightarrow\) not (b) \(\Rightarrow\) not (c) are obvious. If \(a' \neq a\) , there exists by Proposition 1 an ideal \(a''\) of g such that \(a \subset a'' \subset a'\) and dim \(a''/a = 1\) . Then \(a'' = a + Kx\) , hence \([a'', a''] \subset [a, a] + [x, a] = \{0\}\) and hence (a) is false.)
\item
  \begin{enumerate}
  \def\labelenumii{(\alph{enumii})}
  \tightlist
  \item
    Let V be a finite-dimensional vector space over \(K\), \(\mathfrak{g}\) a Lie algebra of nilpotent endomorphisms of V and \((\mathrm{V}_{r})_{0 \leq r \leq n}\) a decreasing sequence of vector subspaces of V with \(V_{0} = V, V_{n} = \{0\}\) and \(\mathfrak{g}(\mathrm{V}_{r}) \subset V_{r+1}\) for \(0 \leq r < n\) . Show by induction on i that \((\mathcal{D}^{i}\mathfrak{g})(\mathrm{V}_{r}) \subset V_{r+2^{i}}\) . If \(\dim V \geq 2^{i}\) , the subspace of elements of V annihilated by \(\mathcal{D}^{i}\mathfrak{g}\) is of dimension \(\geqslant 2^{i}\) . If \(\dim V \leqslant 2^{i}\) , \(\mathcal{D}^{i}\mathfrak{g} = \{0\}\) .
  \end{enumerate}
\end{enumerate}

\begin{enumerate}
\def\labelenumi{(\alph{enumi})}
\setcounter{enumi}{1}
\item
  Let \(\mathfrak{g}\) be a nilpotent Lie algebra and \(i\) an integer \(\geqslant 0\) . If \(\dim \mathcal{D}^{i}\mathfrak{g} > 2^{i} + 1\) the centre of \(\mathcal{D}^{i}\mathfrak{g}\) is of dimension \(\geqslant 2^{i}\) . If \(\dim \mathcal{D}^{i}\mathfrak{g} \leqslant 2^{i} + 1\) , \(\mathcal{D}^{i}\mathfrak{g}\) is commutative. (Apply (a) to the restrictions to \(\mathcal{D}^{i}\mathfrak{g}\) of the \(\operatorname{ad}_{\mathfrak{g}} x, x \in \mathfrak{g}\) .)
\item
  Let \(\mathfrak{g}\) be a nilpotent Lie algebra and \(i\) an integer \(\geqslant 0\) . If \(\mathcal{D}^{i}\mathfrak{g}\) is not commutative, \(\mathcal{D}^{i}\mathfrak{g}/\mathcal{D}^{i+1}\mathfrak{g}\) is of dimension \(\geqslant 2^{i} + 1\) . (Reduce it, by passing to the quotient, to the case where \(\dim \mathcal{D}^{i+1}\mathfrak{g} = 1\) and use (b).)
\end{enumerate}

\begin{enumerate}
\def\labelenumi{\arabic{enumi}.}
\setcounter{enumi}{7}
\tightlist
\item
  Let \(\mathfrak{g}\) be a Lie algebra, \(\mathfrak{m}\) an ideal of codimension 1, \(x\) an element of \(\mathfrak{g}\) not belonging to \(\mathfrak{m}\) and \(z \in \mathfrak{g}\) .
\end{enumerate}

\begin{enumerate}
\def\labelenumi{(\alph{enumi})}
\item
  Show that the linear mapping D which reduces to 0 on m and maps x to z is a derivation if z belongs to the centralizer a of m in g.
\item
  Let \(q\) be the largest integer such that \(a \subset \mathcal{C}^q\mathfrak{g}\) . Show that if further \(z \notin \mathcal{C}^{q+1}\mathfrak{g}\) , \(D\) is not an inner derivation of \(\mathfrak{g}\) .
\item
  Deduce from (a) and (b) and Exercise 7 (c) that if g is a nilpotent Lie algebra of dimension \textgreater1, the vector space of inner derivations of g is of co-dimension \(\geqslant2\) in the space of all derivations of g.
\end{enumerate}

\begin{enumerate}
\def\labelenumi{\arabic{enumi}.}
\setcounter{enumi}{8}
\tightlist
\item
  \begin{enumerate}
  \def\labelenumii{(\alph{enumii})}
  \tightlist
  \item
    Verify that the following multiplication tables define two nilpotent Lie algebras \(\mathfrak{g}_3, \mathfrak{g}_4\) of dimensions 3 and 4:
  \end{enumerate}
\end{enumerate}

\[
\mathfrak {g} _ {3}: \left[ x _ {1}, x _ {2} \right] = x _ {3}, \quad \left[ x _ {1}, x _ {3} \right] = \left[ x _ {2}, x _ {3} \right] = 0
\]

\[
\mathfrak {g} _ {4}: \left[ x _ {1}, x _ {2} \right] = x _ {3}, \quad \left[ x _ {1}, x _ {3} \right] = x _ {4}, \quad \left[ x _ {1}, x _ {4} \right] = \left[ x _ {2}, x _ {3} \right] = \left[ x _ {2}, x _ {4} \right] = \left[ x _ {3}, x _ {4} \right] = 0.
\]

\begin{enumerate}
\def\labelenumi{(\alph{enumi})}
\setcounter{enumi}{1}
\item
  Show that \(\mathfrak{g}_3\) is isomorphic to \(\mathfrak{n}(3, \mathbf{K})\) and also to \(\mathfrak{sl}(2, \mathbf{K})\) if \(\mathbf{K}\) is of characteristic 2.
\item
  Let \(\mathfrak{g}_1\) be the unique Lie algebra of dimension 1. Show that the nilpotent Lie algebras of dimension \(\leqslant 4\) are given by the following table:
\end{enumerate}

\[
\text { dimension } 1: \mathfrak {g} _ {1}; \text { dimension } 2: (\mathfrak {g} _ {1}) ^ {2};
\]

dimension 3: \((\mathfrak{g}_1)^3, \mathfrak{g}_3\) ; dimension 4: \((\mathfrak{g}_1)^4, \mathfrak{g}_3 \times \mathfrak{g}_1, \mathfrak{g}_4\) .

(Use Exercise 7 (c). If \(\mathfrak{g}\) is nilpotent and 3-dimensional and \(\dim \mathcal{D}\mathfrak{g} = 1\) , observe that \(\mathcal{D}\mathfrak{g}\) is contained in the centre of \(\mathfrak{g}\) (Exercise 3 (a)), whence \(\mathfrak{g} = \mathfrak{g}_3\) . If dim \(\mathfrak{g} = 4\) , dim \(\mathcal{D}\mathfrak{g} = 1\) , observe that the bracket on \(\mathfrak{g}\) defines an alternating bilinear form on \(\mathfrak{g} / \mathcal{D}\mathfrak{g}\) , which is necessarily degenerate and hence that there exists a subalgebra \(\mathfrak{h}\) of \(\mathfrak{g}\) contained in the centre of \(\mathfrak{g}\) with dim \(\mathfrak{h} = 1\) , \(\mathfrak{h}\cap \mathcal{D}\mathfrak{g} = \{0\}\) , whence \(\mathfrak{g} = \mathfrak{h}\times \mathfrak{g}_3\) . If dim \(\mathfrak{g} = 4\) , dim \(\mathcal{D}\mathfrak{g} = 2\) , \(\mathcal{D}\mathfrak{g}\) is commutative; applying Theorem 1 to the restrictions to \(\mathcal{D}\mathfrak{g}\) of the \(\operatorname{ad}_{\mathfrak{g}} x\) \((x\in \mathfrak{g})\) , show that there exists a commutative ideal \(\mathfrak{h}\) of \(\mathfrak{g}\) with \(\mathfrak{h}\supseteq \mathcal{D}\mathfrak{g}\) , dim \(\mathfrak{h} = 3\) ; let \(x\in \mathfrak{g},\) \(x\notin \mathfrak{h}\) ; choose a basis of \(\mathfrak{h}\) such that the restriction of \(\operatorname{ad}_{\mathfrak{g}}x\) to \(\mathfrak{h}\) has a Jordan matrix with respect to this basis; then \(\mathfrak{g} = \mathfrak{g}_4\).

\begin{enumerate}
\def\labelenumi{\arabic{enumi}.}
\setcounter{enumi}{9}
\tightlist
\item
  Let \(\mathcal{D}\) be the Lie algebra of derivations of the algebra \(\mathfrak{g}_4\) (Exercise 9). Show that \(\mathcal{D}\) is of dimension 7 and of zero centre and that the ideal \(\mathcal{D}' \subset \mathcal{D}\) of inner derivations has no supplementary subalgebra in \(\mathcal{D}\) .
\end{enumerate}

¶ 11. Let L be a commutative ring, A a left Artinian algebra over L and γ a mapping of A × A into L. Define on A the internal law of composition \((a, b) \mapsto a * b = ab + \gamma(a, b)ba\) . For every subset E of A, let \(\tilde{E}\) denote the subring (with no unit element) of A generated by E. Show that if E consists of nilpotent elements and is stable under the law *, then \(\tilde{E}\) is nilpotent (that is, there exists n \textgreater{} 0 such that \(\tilde{E}^{n} = \{0\}\) ). Proceed as follows:

\begin{enumerate}
\def\labelenumi{\arabic{enumi}.}
\item
  Suppose first that A is a simple ring and hence isomorphic to \(\mathcal{L}_{\mathrm{D}}(\mathrm{T})\) , where T is a left vector space of finite dimension \(m\) over a field D. Argue by induction on \(m\) . Let \(\Phi\) be the set of subsets \(\mathrm{F} \subset \mathrm{E}\) which are stable under \(*\) and such that \(\tilde{\mathrm{F}}\) is nilpotent. Show that \(\tilde{\Phi}\) admits a maximal element M (note that \(\tilde{\mathrm{F}}^m = \{0\}\) for all \(\mathrm{F} \in \Phi\) ). Suppose \(\tilde{\mathbf{M}} \neq \tilde{\mathbf{E}}\) . Show that there exists \(a \in \mathrm{E}\) such that \(a \notin \tilde{\mathbf{M}}\) and \(t * a \in \tilde{\mathbf{M}}\) for all \(t \in \mathbf{M}\) (observe that there can be no infinite sequence \((a_n)\) such that \(a_n \in \mathrm{E}\) , \(a_n \notin \mathbf{M}\) , \(a_n = t_{n-1} * a_{n-1}\) with \(t_{n-1} \in \mathbf{M}\) ). Let S be the subspace of T the direct sum of the \(u(\mathrm{T})\) with \(u \in \tilde{\mathbf{M}}\) ; show that \(S \neq \{0\}\) and \(S \neq T\) and that \(a(S) \subset S\) . Let N be the set of \(v \in E\) such that \(v(S) \subset S\) . By using the induction hypothesis and considering the elements of N as operating on S and on T/S, show that \(N \in \Phi\) , which implies a contradiction.
\item
  In the general case, use the fact that the Jacobson radical of A is nilpotent. Deduce from this result a new proof of Engel's Theorem.
\item
  Let \(\mathfrak{g}\) be a Lie algebra and \(\mathcal{C}^{\infty}\mathfrak{g}\) the intersection of the \(\mathcal{C}^p\mathfrak{g}\) .
\end{enumerate}

\begin{enumerate}
\def\labelenumi{(\alph{enumi})}
\item
  The Lie algebra \(\mathfrak{g} / \mathcal{C}^{\infty}\mathfrak{g}\) is nilpotent.
\item
  Show that there exists a nilpotent subalgebra \(\mathfrak{h}\) of \(\mathfrak{g}\) such that \(\mathfrak{g} = \mathfrak{h} + \mathcal{C}^{\infty}\mathfrak{g}\) . (Argue by induction on \(\dim \mathfrak{g}\) . Suppose that \(\mathfrak{g}\) is not nilpotent. Let \(x\in \mathfrak{g}\) be such that \(\mathbf{I} = \bigcap_{k}\left(\mathrm{(ad}x)^{k}(\mathfrak{g})\neq \{0\}\right)\) . Let \(\mathfrak{n}\) be the union of the kernels of the \((\mathrm{ad}x)^k\) . Show that this is a subalgebra of \(\mathfrak{g}\) and that \(\mathfrak{g}\) is the direct sum of \(\mathbf{I}\) and \(\mathfrak{n}\) . By the induction hypothesis \(\mathfrak{n}\) is the sum of \(\mathcal{C}^{\infty}\mathfrak{n}\) and a nilpotent subalgebra \(\mathfrak{h}\) . Finally, \(\mathcal{C}^{\infty}\mathfrak{n}\subset \mathcal{C}^{\infty}\mathfrak{g}\) and \(\mathbf{I}\subset \mathcal{C}^{\infty}\mathfrak{g}\) .
\item
  Verify that the following multiplication table defines a (solvable) Lie algebra g of dimension 5:
\end{enumerate}

\[
\left[ x _ {1}, x _ {2} \right] = x _ {5}, \quad \left[ x _ {1}, x _ {3} \right] = x _ {3}, \quad \left[ x _ {1}, x _ {4} \right] = - x _ {4}, \quad \left[ x _ {3}, x _ {4} \right] = x _ {5}
\]

and the other \([x_i, x_j]\) are zero. Show that for this Lie algebra

\[
\mathcal {C} ^ {\infty} \mathfrak {g} = \mathrm{K} x _ {3} + \mathrm{K} x _ {4} + \mathrm{K} x _ {5}
\]

but that there exists no supplementary subalgebra of \(\mathcal{C}^{\infty}\mathfrak{g}\) in \(\mathfrak{g}\) .

¶ 13. Let g be a Lie algebra and \(\mathfrak{z}\) the centralizer of \(C^{\infty}g\) in g. If \(\mathfrak{z} \not\subset C^{\infty}g\) , g has non-zero centre. (Write \(g = C^{\infty}g + t\) , where t is a nilpotent subalgebra of g (Exercise 12)). Let \(g_{1} = \mathfrak{z} + t\) . Let b be a nilpotent subalgebra such that \(g_{1} = C^{\infty}g_{1} + \mathfrak{b}\) . Show that \(g = C^{\infty}g + \mathfrak{b}\) and that \(C^{\infty}g_{1} \subset \mathfrak{z} \cap C^{\infty}g\) . Let y be an element of \(\mathfrak{z}\) not belonging to \(C^{\infty}g\) . Write \(y = x + x'\) where \(x \in \mathfrak{b}, x' \in C^{\infty}g_{1}\) ; then \(x \in \mathfrak{z}\) and \(x \neq 0\) and hence \(\mathfrak{z} \cap \mathfrak{b}\) is a non-zero ideal of b; using Exercise 3 (a), deduce that there exists a non-zero element of g permutable with \(C^{\infty}g\) and b.)

¶ 14. Let g be a Lie algebra and b a subinvariant subalgebra of g such that the centralizer \(\mathfrak{z}(b)\) of b in g is zero.

\begin{enumerate}
\def\labelenumi{(\alph{enumi})}
\item
  The centralizer \(\mathfrak{z}(\mathcal{C}^{\infty}\mathfrak{b})\) of \(\mathcal{C}^{\infty}\mathfrak{b}\) in \(\mathfrak{g}\) is contained in \(\mathcal{C}^{\infty}\mathfrak{b}\) . (If \(\mathfrak{z}(\mathcal{C}^{\infty}\mathfrak{b})\not\subset\mathcal{C}^{\infty}\mathfrak{b},\mathfrak{z}(\mathcal{C}^{\infty}\mathfrak{b})\not\subset\mathfrak{b}\) by Exercise 13; \(\mathcal{C}^{\infty}\mathfrak{b}\) is an ideal of \(\mathfrak{g}\) (§ 1, Exercise 14); let \(\mathfrak{c}=\mathfrak{b}+\mathfrak{z}(\mathcal{C}^{\infty}\mathfrak{b})\) ; \(\mathfrak{c}\) is a subalgebra, \(\mathfrak{b}\) is subinvariant in \(\mathfrak{c}\) ; \(\mathfrak{b}\) is an ideal of \(\mathfrak{b}_{1}\subset\mathfrak{c}\) , where \(\mathfrak{b}\neq\mathfrak{b}_{1}\) ; let \(y\in\mathfrak{b}_{1},y\notin\mathfrak{b};y=x+x'\) , where \(x\in\mathfrak{z}(\mathcal{C}^{\infty}\mathfrak{b}),x'\in\mathfrak{b}\) ; then \(x\in\mathfrak{z}(\mathcal{C}^{\infty}\mathfrak{b})\cap\mathfrak{b}_{1},x\notin\mathfrak{b};\mathfrak{b}_{1}=\mathrm{Kx}+\mathfrak{b}\) is a subalgebra; \(\mathcal{C}^{\infty}\mathfrak{c}=\mathcal{C}^{\infty}\mathfrak{b}\) because \(\mathfrak{c}\cap\mathfrak{z}(\mathcal{C}^{\infty}\mathfrak{b})\subset\mathcal{C}^{\infty}\mathfrak{b}\) , whence \(\mathcal{C}^{k}\mathfrak{c}\subset\mathcal{C}^{k}\mathfrak{b}\) for all \(k\) ; \(\mathfrak{z}(\mathcal{C}^{\infty}\mathfrak{b})\cap\mathfrak{b}\not\subset\mathcal{C}^{\infty}\mathfrak{b}\) and hence the centre of \(\mathfrak{b}\) is non-zero by Exercise 13; a fortiori, \(\mathfrak{z}(\mathfrak{b})\neq\{0\}\) , which is absurd.)
\item
  Deduce from (a) that, if \(\mathfrak{D}\) denotes the Lie algebra of derivations of \(\mathcal{C}^{\infty}\mathfrak{b}\) and \(\mathfrak{a}\) the centre of \(\mathcal{C}^{\infty}\mathfrak{b}\) , then \(\dim \mathfrak{g} \leqslant \dim \mathfrak{D} + \dim \mathfrak{a}\) (let \(\mathfrak{g}\) operate on \(\mathcal{C}^{\infty}\mathfrak{b}\) by the adjoint representation).
\end{enumerate}

\begin{enumerate}
\def\labelenumi{\arabic{enumi}.}
\setcounter{enumi}{14}
\tightlist
\item
  Let \(l\) be a Lie algebra with zero centre, \(\mathfrak{D}_1\) the Lie algebra of derivations of \(l\) , \(\mathfrak{D}_2\) the Lie algebra of derivations of \(\mathfrak{D}_1, \ldots\) . Then \(l\) is an ideal of \(\mathfrak{D}_1, \mathfrak{D}_1\) is an ideal of \(\mathfrak{D}_2\) , etc. (§ 1, Exercise 15 (c)). Let \(\mathfrak{D}\) be the Lie algebra of derivations of \(\mathcal{C}^\infty l\) and \(a\) the centre of \(\mathcal{C}^\infty l\) .
\end{enumerate}

\begin{enumerate}
\def\labelenumi{(\alph{enumi})}
\item
  Then \(\dim \mathfrak{D}_1 \leqslant \dim \mathfrak{D} + \dim \mathfrak{a}\) . (Use Exercise 14 above and Exercise 15 of § 1.)
\item
  Deduce from (a) that for i sufficiently large all the derivations of \(D_{i}\) are inner.
\end{enumerate}

\begin{enumerate}
\def\labelenumi{\arabic{enumi}.}
\setcounter{enumi}{15}
\item
  Let \(\mathfrak{g}_1\) and \(\mathfrak{g}_2\) be Lie K-algebras and \(\mathfrak{n}_1\) and \(\mathfrak{n}_2\) their largest nilpotent ideals. Show that the largest nilpotent ideal of \(\mathfrak{g}_1 \times \mathfrak{g}_2\) is \(\mathfrak{n}_1 \times \mathfrak{n}_2\) .
\item
  We reconsider the Lie algebra \(\mathfrak{g}_3\) of Exercise 9 (a) with its basis \((x_1, x_2, x_3)\) .
\end{enumerate}

\begin{enumerate}
\def\labelenumi{(\alph{enumi})}
\tightlist
\item
  Let V be the vector space K{[}X{]}. Let D be the differential operator with respect to X on V and M the operator of multiplication by X on V. Show that if K is of characteristic 0 the mapping
\end{enumerate}

\[
\alpha x _ {1} + \beta x _ {2} + \gamma x _ {3} \mapsto \alpha D + \beta M + \gamma
\]

is an infinite-dimensional simple representation \(\rho\) of \(\mathfrak{g}\) on V.

\begin{enumerate}
\def\labelenumi{(\alph{enumi})}
\setcounter{enumi}{1}
\tightlist
\item
  If K is of characteristic p \textgreater{} 0, the ideal \((\mathbf{X}^{p})\) of K{[}X{]} is stable under \(\rho(\mathfrak{g})\) . Under the quotient representation of g on \(\mathrm{K}[\mathrm{X}] / (\mathrm{X}^{p})\) , no line is stable.
\end{enumerate}

\begin{enumerate}
\def\labelenumi{\arabic{enumi}.}
\setcounter{enumi}{17}
\tightlist
\item
  Let \(\mathfrak{g}\) be a 7-dimensional vector space over \(\mathbf{K}\) with a basis \((e_i)_{1\leqslant i\leqslant 7}\) . An alternating bracket is defined on \(\mathfrak{g}\) by the formulae:
\end{enumerate}

\[
\left[ e _ {i}, e _ {j} \right] = \alpha_ {i j} e _ {i + j} \quad (1 \leqslant i <   j \leqslant 7, i + j \leqslant 7)\tag{1}
\]

and all the other brackets \([e_i, e_j]\) are zero for \(i < j\) .

\begin{enumerate}
\def\labelenumi{(\alph{enumi})}
\tightlist
\item
  For the Jacobi identity to hold, it is necessary and sufficient that:
\end{enumerate}

\begin{enumerate}
\def\labelenumi{(\arabic{enumi})}
\setcounter{enumi}{1}
\tightlist
\item
\end{enumerate}

\[
- \alpha_ {2 3} \alpha_ {1 5} + \alpha_ {1 3} \alpha_ {2 4} = 0\tag{2}
\]

\[
\alpha_ {1 2} \alpha_ {3 4} - \alpha_ {2 4} \alpha_ {1 6} + \alpha_ {1 4} \alpha_ {2 5} = 0.\tag{3}
\]

\begin{enumerate}
\def\labelenumi{(\alph{enumi})}
\setcounter{enumi}{1}
\item
  Suppose henceforth that all the \(\alpha_{ij}\) are \(\neq 0\) . Show that the ideals \(\mathcal{C}^2\mathfrak{g}\) , \(\mathcal{C}^3\mathfrak{g}\) , \(\mathcal{C}^4\mathfrak{g}\) , \(\mathcal{C}^5\mathfrak{g}\) , \(\mathcal{C}^6\mathfrak{g}\) admit the following bases: \((e_3, e_4, e_5, e_6, e_7)\) , \((e_4, e_5, e_6, e_7)\) , \((e_5, e_6, e_7)\) , \((e_6, e_7)\) , \((e_7)\) . Show that the centralizer \(\mathfrak{h}\) of \(\mathcal{C}^5\mathfrak{g}\) admits the basis \((e_2, e_3, e_4, e_5, e_6, e_7)\) .
\item
  Let \((e_i')_{1 \leqslant i \leqslant 7}\) be another basis of \(\mathfrak{g}\) such that \(\mathcal{C}^6\mathfrak{g}\) admits the basis \((e_7')\) , \(\mathcal{C}^5\mathfrak{g}\) admits the basis \((e_6', e_7')\) , \ldots, \(\mathfrak{h}\) admits the basis \((e_2', e_3', e_4', e_5', e_6', e_7')\) . Show that
\end{enumerate}

\[
\left[ e _ {i} ^ {\prime}, e _ {j} ^ {\prime} \right] = \alpha_ {i j} ^ {\prime} e _ {i + j} ^ {\prime} + \beta e _ {i + j + 1} ^ {\prime} + \gamma e _ {i + j + 2} ^ {\prime} + \dots + \lambda e _ {7} ^ {\prime},
\]

where

\[
\alpha_ {1 4} ^ {\prime} \alpha_ {2 5} ^ {\prime} \alpha_ {1 6} ^ {\prime - 1} \alpha_ {2 4} ^ {\prime - 1} = \alpha_ {1 4} \alpha_ {2 5} \alpha_ {1 6} ^ {- 1} \alpha_ {2 4} ^ {- 1}.
\]

\begin{enumerate}
\def\labelenumi{(\alph{enumi})}
\setcounter{enumi}{3}
\tightlist
\item
  Deduce that there exists, if K is infinite, an infinity of non-isomorphic nilpotent Lie algebras of dimension 7 over K and that there exist complex nilpotent Lie algebras which cannot be derived from real nilpotent Lie algebras by extending the base field from R to C.
\end{enumerate}

\begin{enumerate}
\def\labelenumi{\arabic{enumi}.}
\setcounter{enumi}{18}
\tightlist
\item
  \begin{enumerate}
  \def\labelenumii{(\alph{enumii})}
  \tightlist
  \item
    Let \(\mathfrak{g}\) be a Lie algebra. Let \(\mathfrak{S}\) be the Lie algebra of derivations of \(\mathfrak{g}\) . Characteristic ideals \(\mathfrak{g}^{[k]}\) of \(\mathfrak{g}\) are defined inductively as follows: \(\mathfrak{g}^{[0]} = \mathfrak{g}\) and \(\mathfrak{g}^{[k+1]}\) is the subspace of \(\mathfrak{g}\) generated by the Dx ( \(\mathrm{D} \in \mathfrak{S}, x \in \mathfrak{g}^{[k]}\) ). The following conditions are equivalent: (1) every derivation of \(\mathfrak{g}\) is nilpotent: (2) \(\mathfrak{g}^{[k]} = \{0\}\) for \(k\) sufficiently large; (3) the holomorph of \(\mathfrak{g}\) is nilpotent. If they are fulfilled, \(\mathfrak{g}\) is said to be characteristically nilpotent. Such an algebra is nilpotent.
  \end{enumerate}
\end{enumerate}

\begin{enumerate}
\def\labelenumi{(\alph{enumi})}
\setcounter{enumi}{1}
\tightlist
\item
  Verify that an 8-dimensional Lie algebra is defined by the following multiplication table:
\end{enumerate}

\[
\begin{array}{r} \left[ x _ {1}, x _ {2} \right] = x _ {3}, \quad \left[ x _ {1}, x _ {3} \right] = x _ {4}, \quad \left[ x _ {1}, x _ {4} \right] = x _ {5}, \quad \left[ x _ {1}, x _ {5} \right] = x _ {6} \\ \left[ x _ {1}, x _ {6} \right] = x _ {8}, \quad \left[ x _ {1}, x _ {7} \right] = x _ {6}, \quad \left[ x _ {2}, x _ {3} \right] = x _ {5}, \quad \left[ x _ {2}, x _ {4} \right] = x _ {6} \\ \left[ x _ {2}, x _ {5} \right] = x _ {7}, \quad \left[ x _ {2}, x _ {6} \right] = 2 x _ {8}, \quad \left[ x _ {3}, x _ {4} \right] = - x _ {7} + x _ {8}, \quad \left[ x _ {3}, x _ {5} \right] = - x _ {8} \end{array}
\]

and \([x_i, x_j] = 0\) for \(i + j > 8\) . Show that

\[
\begin{array}{c c} \mathcal {C} ^ {2} \mathfrak {g} = \sum_ {i = 3} ^ {8} \mathrm{K} x _ {i}, & \mathcal {C} ^ {3} \mathfrak {g} = \sum_ {i = 4} ^ {8} \mathrm{K} x _ {i}, \\ & \mathcal {C} ^ {6} \mathfrak {g} = \mathrm{K} x _ {8}, \quad \mathcal {C} ^ {7} \mathfrak {g} = \{0 \}, \end{array} \quad \begin{array}{c c} \mathcal {C} ^ {4} \mathfrak {g} = \sum_ {i = 5} ^ {8} \mathrm{K} x _ {i}, & \mathcal {C} ^ {5} \mathfrak {g} = \sum_ {i = 6} ^ {8} \mathrm{K} x _ {i}, \\ & [ \mathcal {C} ^ {2} \mathfrak {g}, \mathcal {C} ^ {2} \mathfrak {g} ] = \mathrm{K} x _ {7} + \mathrm{K} x _ {8}, \end{array}
\]

and that the transporter of \(\mathcal{C}^2\mathfrak{g}\) into \(\mathcal{C}^4\mathfrak{g}\) is \(\sum_{i=2}^{\infty}\mathrm{Kx}_i\) . Deduce that every derivation D of \(\mathfrak{g}\) is defined by formulae \(\mathrm{D}x_i = \sum_{j=1}^{8}u_{ij}x_j\) . Then show that \(u_{ii} = 0\) for all \(i\) if the characteristic of K differs from 2, so that \(\mathfrak{g}\) is characteristically nilpotent.

\begin{enumerate}
\def\labelenumi{(\alph{enumi})}
\setcounter{enumi}{2}
\tightlist
\item
  If the characteristic of K is \(\neq2\) , show that the Lie algebra g of (b) is not the derived Lie algebra of any Lie algebra. (If g = Dh, show first that h is nilpotent. Observe that \(\dim g/Dg = 2\) , which with Exercise 7 (c) leads to a contradiction.)
\end{enumerate}

¶ 20. Let g be a Lie algebra and U its enveloping algebra.

\begin{enumerate}
\def\labelenumi{(\alph{enumi})}
\item
  Let \(\mathfrak{g}'\) be an ideal of \(\mathfrak{g}\) , \(\mathrm{U}' \subset \mathrm{U}\) its enveloping algebra, \(x \in \mathfrak{g}\) and \(a_1, \ldots, a_n\) in U. Suppose that \(a_i = x\) for the indices \(j_1, \ldots, j_n\) and \(a_i \in \mathrm{U}'\) for the other indices (let \(k_1, \ldots, k_q\) be these indices, with \(k_1 < k_2 < \cdots < k_q\) ). Then \(a_1 a_2 \ldots a_n - x^p a_{k_1}a_{k_2} \ldots a_{k_q}\) is a sum of terms of the form \(x^p b\) , where \(b \in \mathrm{U}'\) and \(p' < p\) . (Argue by induction on \(p\) .)
\item
  Suppose henceforth that g is nilpotent and K is of characteristic 0. Let g' be an ideal of codimension 1 in g, U' its enveloping algebra, x an element of g not belonging to g' and U₀ (resp. U₀') the subalgebra of U (resp. U') annihilated by a set b of derivations of g (which extend to derivations of U) mapping g into g'. Suppose that U₀ is contained in the centre of U and U₀ ≠ U'. Show that there exist a₁ ∈ U₀', a₂ ∈ U' such that a₁ ≠ 0 and a = xa₁ + a₂ ∈ U₀. (Let xᵐbₘ + xᵐ⁻¹bₘ₋₁ + ⋯ + b₀ ∈ U₀ with bₘ, \ldots, b₀ in U', m \textgreater{} 0, bₘ ≠ 0. Show, using ⟨a⟩, that, for every derivation D ∈ b of g, Dbₘ = 0, m(Dx)bₘ + Dbₘ₋₁ = 0 and hence D(mxbₘ + bₘ₋₁) = 0.) Show that U₀ is contained in the algebra K{[}a, a₁⁻¹, U₀'{]} generated by a, a₁⁻¹, U₀' in the field of fractions of U₀ (which can be formed by Corollary 7 to Theorem 1 of § 2). Deduce that the field of fractions of U₀ is the field generated by a and U₀'. Show that a is transcendental over K(U₀').
\item
  Let \(\{0\} = \mathfrak{g}_0 \subset \mathfrak{g}_1 \subset \cdots \subset \mathfrak{g}_n = \mathfrak{g}\) be a sequence of ideals of \(\mathfrak{g}\) of dimensions \(0, 1, \ldots, n\) . Let \(x_j\) be an element of \(\mathfrak{g}_j\) not belonging to \(\mathfrak{g}_{j-1}\) . Let \(j_1 < j_2 < \cdots < j_n\) be the indices \(j\) such that there exists in \(\mathrm{U}^j\) (enveloping algebra of \(\mathfrak{g}_j\) ) an element of the centre \(\mathrm{U}_0\) of \(\mathrm{U}\) not belonging to \(\mathrm{U}^{j-1}\) . By (b) there exists \(a_{j_1} \in \mathrm{U}_0 \cap \mathrm{U}^{j_1}\) and \(a_{j_2} \in \mathrm{U}^{j_2-1}\) such that \(a_{j_1} \neq 0\) and \(a_j = x_j a_{j_1} + a_{j_2} \in \mathrm{U}_0 \cap \mathrm{U}^{j}\) . Show by induction on \(n\) that
\end{enumerate}

\[
\mathrm{U} _ {0} \subset \mathrm{K} [ a _ {j _ {1}}, \dots , a _ {j _ {q}}, a _ {j _ {1}} ^ {- 1}, \dots , a _ {j _ {q}} ^ {- 1} ]
\]

and that the field of fractions of \(U_{0}\) is generated by the algebraically independent elements \(a_{j_{1}},\ldots,a_{j_{q}}\) . In particular, the field of fractions of the centre of U is a pure transcendental extension of K.

\begin{enumerate}
\def\labelenumi{\arabic{enumi}.}
\setcounter{enumi}{20}
\tightlist
\item
  Let g be a Lie algebra and σ an automorphism of g.
\end{enumerate}

\begin{enumerate}
\def\labelenumi{(\alph{enumi})}
\tightlist
\item
  Suppose first that \(\mathbf{K}\) is algebraically closed; for all \(\lambda \in \mathbf{K}\) let \(\mathrm{L}_{\lambda}\) be the set of \(x \in \mathfrak{g}\) which are annihilated by a power of \(\sigma - \lambda \mathrm{I}\) ; \(\mathfrak{g}\) is the direct sum of the \(\mathrm{L}_{\lambda}\) . Show that \([\mathrm{L}_{\lambda}, \mathrm{L}_{\mu}] \subset \mathrm{L}_{\lambda \mu}\) . (Note that
\end{enumerate}

\[
(\sigma - \lambda \mu \mathrm{I}) ([ x, y ]) = [ (\sigma - \lambda \mathrm{I}) x, \sigma y ] + [ \lambda x, (\sigma - \mu \mathrm{I}) y ].)
\]

\begin{enumerate}
\def\labelenumi{(\alph{enumi})}
\setcounter{enumi}{1}
\item
  With an arbitrary field K, suppose that none of the eigenvalues of \(\sigma\) (in an algebraically closed extension of K) is a root of unity. Show that g is nilpotent. (Reduce it to the case where K is algebraically closed. If \(\lambda_{1},\ldots,\lambda_{m}\) are the distinct eigenvalues of \(\sigma\) , the \(\lambda_{i}\lambda_{j},\lambda_{i}\lambda_{j}^{2},\lambda_{i}\lambda_{j}^{3},\ldots\) are not all eigenvalues and hence \(ad_{g}x\) is nilpotent for \(x\in L_{\lambda_{i}}\) . Conclude using Exercise 11 applied to the set of \(ad_{g}x\) , where x runs through the union of the \(L_{\lambda_{i}}\) ).
\item
  With an arbitrary field K, suppose that \(\sigma^{q}=I\) , where q is a prime number and that no eigenvalue of \(\sigma\) is equal to 1. Show that g is nilpotent. (Same method as in (b), observing that, if \(g\neq\{0\}\) , q is not equal to the characteristic of K and that for every ordered pair of eigenvalues \(\lambda_{i},\lambda_{j}\) of \(\sigma\) there exists an integer k such that \(\lambda_{i}\lambda_{j}^{k}=1\) .)
\end{enumerate}

\begin{enumerate}
\def\labelenumi{\arabic{enumi}.}
\setcounter{enumi}{21}
\tightlist
\item
  \begin{enumerate}
  \def\labelenumii{(\alph{enumii})}
  \tightlist
  \item
    Let \(u, v\) be two endomorphisms of a vector space \(\mathbf{E}\) of finite dimension over an algebraically closed field \(\mathbf{K}\) . For every eigenvalue \(\lambda\) of \(u\) let \(\mathrm{E}_{\lambda}\) be the subspace of \(\mathbf{E}\) consisting of the vectors annihilated by a power of \(u - \lambda \mathrm{I}\) . Show that, if \((\operatorname{ad} u)^n v = 0\) , the subspaces \(\mathrm{E}_{\lambda}\) are stable under \(v\) .
  \end{enumerate}
\end{enumerate}

\begin{enumerate}
\def\labelenumi{(\alph{enumi})}
\setcounter{enumi}{1}
\item
  Deduce from (a) that, if g is a nilpotent Lie algebra of endomorphisms of E, E is the direct sum of subspaces \(F_{j}\) ( \(1 \leqslant j \leqslant m\) ) which are stable under g and such that on each \(F_{j}\) the restriction of every element \(u \in g\) can be written \(\lambda_{j}(u)\mathbf{I} + u_{j}\) , where \(\lambda_{j}(u) \in \mathbf{K}\) and \(u_{j}\) is nilpotent.
\item
  If \(\mathbf{K}\) is of characteristic 2, \(\mathrm{E} = \mathrm{K}^2\) and \(\mathfrak{g}\) is nilpotent algebra \(\mathfrak{sl}(2,\mathbf{K})\) , show that \(m = 1\) and that possibly \(\lambda (u + v)\neq \lambda (u) + \lambda (v)\) for two elements \(u,v\) of \(\mathfrak{g}\) . (For a sequel to this exercise, cf.~§ 5, Exercise 12.)
\end{enumerate}

\begin{enumerate}
\def\labelenumi{\arabic{enumi}.}
\setcounter{enumi}{22}
\tightlist
\item
  Let \(\mathfrak{g}\) be a Lie \(p\) -algebra over a perfect field \(\mathbf{K}\) of characteristic \(p > 0\) . \(\mathfrak{g}\) is called \(p\) -unipotent if for all \(x \in \mathfrak{g}\) there exists \(m\) such that \(x^{p^m} = 0\) .
\end{enumerate}

\begin{enumerate}
\def\labelenumi{(\alph{enumi})}
\item
  Show that every \(p\) -unipotent Lie \(p\) -algebra is nilpotent.
\item
  Suppose that \(\mathfrak{g}\) is nilpotent and let \(\mathfrak{h}\) denote the \(p\) -core of the centre of \(\mathfrak{g}\) (§ 1, Exercise 23 (a)). Show that \(\mathfrak{g} / \mathfrak{h}\) is \(p\) -unipotent.
\item
  Let \(\mathfrak{g}\) be the nilpotent Lie \(p\) -algebra with a basis of three elements \(e_1, e_2, e_3\) such that \([e_1, e_2] = [e_1, e_3] = 0\) , \([e_2, e_3] = e_1\) , \(e_1^p = e_1\) , \(e_2^p = e_3^p = 0\) . Then \(\mathfrak{h} = \mathrm{Ke}_1\) , but \(\mathfrak{g}\) is not the direct sum of \(\mathfrak{h}\) and a \(p\) -unipotent \(p\) -subalgebra.
\item
  If g is p-unipotent, show that in the restricted enveloping algebra of g the two-sided ideal generated by g is nilpotent. (Argue by induction on the dimension of g.)
\end{enumerate}

\begin{enumerate}
\def\labelenumi{\arabic{enumi}.}
\setcounter{enumi}{23}
\item
  Suppose that the field K is of characteristic 2. Show that in the nilpotent Lie algebra \(\mathfrak{g}_4\) of Exercise 9 there exists no 2-mapping.
\item
  Let \(\mathfrak{g}\) be a Lie \(p\) -algebra. Show that the largest nilpotent ideal of \(\mathfrak{g}\) is a \(p\) -ideal (cf.~§ 1, Exercise 22).
\item
  Let \(\mathbf{G}\) be a group and let \((\mathbf{H}_n)_{n\geqslant 1}\) be a decreasing sequence of subgroups of \(\mathbf{G}\) ; suppose that \(\mathbf{H}_1 = \mathbf{G}\) and that if we write \((x,y) = xyx^{-1}y^{-1}\) , the relations \(x\in \mathrm{H}_i,y\in \mathrm{H}_j\) imply \((x,y)\in \mathrm{H}_{i + j}\) .
\end{enumerate}

\begin{enumerate}
\def\labelenumi{(\alph{enumi})}
\item
  Let \(G_{i} = H_{i}/H_{i+1}\) ; show that \(G_{i}\) is commutative and that the mapping \(x, y \mapsto (x, y)\) defines on passing to the quotient a Z-bilinear mapping of \(G_{i} \times G_{j}\) into \(G_{i+j}\) .
\item
  Let \(\mathrm{gr}(\mathbf{G}) = \sum_{i=1}^{\infty} \mathbf{G}_i\) and extend by linearity the mappings \(\mathbf{G}_i \times \mathbf{G}_j \to \mathbf{G}_{i+j}\) defined in (a) to a Z-bilinear mapping of \(\mathrm{gr}(\mathbf{G}) \times \mathrm{gr}(\mathbf{G})\) into \(\mathrm{gr}(\mathbf{G})\) . Show that \(\mathrm{gr}(\mathbf{G})\) is a Lie Z-algebra with this mapping (to verify the Jacobi identity, use the following formula:
\end{enumerate}

\[
((x, y), z ^ {y}) \cdot ((y, z), x ^ {z}) \cdot ((z, x), y ^ {x}) = e \quad x, y, z \text {   in   } G,
\]

where \(x^{y}\) denotes \(yxy^{-1}\) and \(e\) denotes the identity element of G).

\begin{enumerate}
\def\labelenumi{(\alph{enumi})}
\setcounter{enumi}{2}
\tightlist
\item
  Suppose that there exists \(n\) such that \(\mathrm{H}_n = \{e\}\) . Show that \(\operatorname{gr}(G)\) is a nilpotent Lie \(\mathbf{Z}\) -algebra.
\end{enumerate}

\begin{enumerate}
\def\labelenumi{\arabic{enumi}.}
\setcounter{enumi}{26}
\tightlist
\item
   Let A be an associative algebra with unit element 1 and let \(\mathrm{A}_0 = \mathrm{A} \supset \mathrm{A}_1 \supset \dots \supset \mathrm{A}_n \supset \dots\) be a decreasing sequence of two-sided ideals of A such that \(\mathrm{A}_i \cdot \mathrm{A}_j \subset \mathrm{A}_{i+j}\) . Let G be a group with identity element \(e\) and let \(f: \mathrm{G} \to \mathrm{A}\) be a mapping such that \(f(e) = 1\) , \(f(xy) = f(x) \cdot f(y)\) and \(1 - f(x) \in \mathrm{A}_1\) for all \(x \in \mathrm{G}\) . Let \(\mathrm{H}_n\) denote the set of \(x \in \mathrm{G}\) such that \(1 - f(x) \in \mathrm{A}_n\) . Show that the \(\mathrm{H}_n\) satisfy the conditions of Exercise 26. Show that the mapping \(x \mapsto f(x) - 1\) defines on taking quotients an injective homomorphism of the Lie algebra \(\mathrm{gr}(\mathrm{G})\) into the Lie algebra associated with the graded ring \(\mathrm{gr}(\mathrm{A}) = \sum \mathrm{A}_n / \mathrm{A}_{n+1}\) (cf.~Commutative Algebra, Chapter III).
\end{enumerate}

% semantic-fragments: legacy-input-seam
\relax{}
